\todo{Lecture 18}
\begin{definition}
Soient $r,g:D\to \mathbb{R}$ et $x_{0}$ adherent a $D\setminus \{ x_{0} \}$. Supposons que $g(x)\neq 0$ pour $x \in D$ suffisament proche de $x_0$, avec $x\neq x_{0}$. On ecrit
$$
r(x)=o(g(x)),\quad x\to x_{0},
$$
si $|r(x)|$ est negligeable compare a $|g(x)|$ pour $x$ proche de $x_{0}$, c-a-d
$$
\lim_{ x \to x_{0},x\neq x_{0} } \frac{r(x)}{g(x)}=0. 
$$
\end{definition}
\begin{definition}
Soient $r,g:D\to \mathbb{R}$ et $x_{0}$ adherent a $D\setminus \{ x_{0} \}$. On ecrit 
$$
r(x)=O(g(x)),\quad x\to x_{0},
$$
s'il existe $C>0$ et $\delta>0$ tels que, pour tout $x \in D$,
$$
0<|x-x_{0}|<\delta \Rightarrow |r(x)|\le C|g(x)|.
$$
\end{definition}

\begin{definition}
Soient $f:D\to \mathbb{R}$, $n\in \mathbb{N}$ et $x_{0}$ adherent a $D\setminus \{ x_{0} \}$. On dit que $f$ admet un developpement limite a l'ordre $n$ en $x_{0}$, ou $\mathrm{DL_{n}}(x_{0})$, s'il existe un polynome $p$ de degre au plus $n$ tel que
$$
f(x)=p(x)+o((x-x_{0})^{n}),\quad x\to x_{0}.
$$
Le polynome $p$ est appele la partie principale du DL.
\end{definition}
\begin{remark}
\begin{itemize}
\item Si $f$ admet un $\mathrm{DL}_{n}(x_{0})$, sa partie principale est unique.
\item Taylor-Young, sous certaines hypotheses, fournit l'existence du $\mathrm{DL}_{n}(x_{0})$ et dit que sa partie principale est le polynome de Taylor a l'ordre $n$.
\item Cela dit, la definition dur $\mathrm{DL}_{n}(x_{0})$ ne requiert pas de derivabilite.
\end{itemize}
\end{remark}